
sur \(x\), d'où le
résultat, l'homomorphisme étant compatible aux filtrations.

\section*{6. Anneau de Chern}

Dans ce paragraphe, l'anneau de base sera un anneau binômial (2.7) \(K\), fixé une fois pour toutes. Nous utiliserons la catégorie \(\mathcal C\) (1.1), et par suite reviendrons aux notations du paragraphe 1, réservant la notation \(\widehat G\) au foncteur défini sur \(\mathcal C\) (1.1), et notant \(\widehat G\circ i\) sa restriction à \(\mathcal C_0\) (1.8).

6.1. Nous avons défini en 3.5 et 3.6 des morphismes de foncteurs :
\[
  (\widehat G\circ i)\times(\widehat G\circ i)\xrightarrow{\nabla}\widehat G\circ i,
  \qquad
  \widehat G^t\circ i\xrightarrow{\gamma'}(\widehat G^u\circ i)\circ(\widehat G^t\circ i)_\nabla,
  \qquad
  \widehat G\circ i\xrightarrow{a}\widehat G\circ i,
\]
le dernier étant la multiplication pour la loi \(\nabla\) par \(\gamma_t(a)\), avec \(a\in K\).

Par 1.9, on en déduit des morphismes :
\[
  \widehat G\times\widehat G\xrightarrow{\nabla}\widehat G_g,
  \qquad
  \widehat G^t\xrightarrow{\gamma'}\widehat G^u\circ i\circ\widehat G^t_{g\nabla},
  \qquad
  \widehat G\xrightarrow{a}\widehat G_g;
\]
et il résulte des propriétés d'isobarité des polynômes \(P'_n\), \(Q'_{i,j}\), \(R_{n,a}\) (3.5, 3.6, 2.9) que l'image du premier et du troisième foncteurs est dans \(\widehat G\), et celle du deuxième, dans le sous-groupe des séries formelles en \(u\) dont les coefficients en degré non nul sont dans \(\widehat G^t\).

En composant le premier et le troisième avec le morphisme canonique : \(\widehat G\longrightarrow 1+\widehat{\operatorname{Gr}\,\mathord{.}\,}^{+}\) (1.11), on obtient d'après 1.12 deux morphismes de
foncteurs sur \(\mathcal G\) :
